\section{Finite similarity and the oblique-projector no-go}
\label{sec:similarity}

Let \(J_N\subset I(P_N;\C)\) be the ideal of incidence functions with zero
diagonal. It is nilpotent, and
\begin{equation}
I(P_N;\C)/J_N\cong\C^{P_N}.
\label{eq:incidence-semisimple}
\end{equation}
The quotient records the coordinate characters. Off-diagonal incidence data
live in the radical.

\begin{theorem}[Uniqueness of complete primitive lifts]
\label{thm:finite-conjugacy}
Let \(\{e_x:x\in P_N\}\) and \(\{f_x:x\in P_N\}\) be complete orthogonal
primitive-idempotent families satisfying
\[
f_x\equiv e_x\pmod{J_N}
\]
for every \(x\). Then there is a unit \(v\in1+J_N\) such that
\begin{equation}
f_x=ve_xv^{-1}\qquad(x\in P_N).
\label{eq:lift-conjugacy}
\end{equation}
\end{theorem}

\begin{proof}
Set
\[
v=\sum_x f_xe_x.
\]
Modulo \(J_N\), one has
\[
v\equiv\sum_x e_x=1.
\]
Thus \(v\) is invertible because \(J_N\) is nilpotent. Orthogonality and
completeness give
\[
f_xv=f_xe_x=ve_x.
\]
Right multiplication by \(v^{-1}\) proves \eqref{eq:lift-conjugacy}.
\end{proof}

For the canonical family, the conclusion is visible without classification:
\begin{equation}
q_n=\zeta\epscoord_n\zeta^{-1}.
\label{eq:canonical-conjugacy}
\end{equation}
Because \(\zeta-\delta\in J_N\), this is exactly a unitriangular change of
basis. If labels are retained, \cref{thm:finite-conjugacy} exhausts the
freedom in a complete lift. If labels are forgotten, a compatible relabeling
may additionally permute the coordinate characters.

\begin{corollary}[Power-trace and determinant collapse]
\label{cor:finite-determinant}
For arbitrary scalars \(b_n\) and every \(r\ge1\),
\begin{equation}
\Tr\left(\sum_n b_nq_n\right)^r=\sum_n b_n^r,
\label{eq:finite-power-trace}
\end{equation}
and
\begin{equation}
\det\left(I-z\sum_n b_nq_n\right)=\prod_n(1-zb_n).
\label{eq:finite-determinant}
\end{equation}
\end{corollary}

\begin{proof}
Either use pair annihilation and \(\Tr q_n=1\), or conjugate the operator to
\(\sum_n b_n\epscoord_n\) using \eqref{eq:canonical-conjugacy}.
\end{proof}

\begin{figure}[t]
\centering
\resizebox{0.94\textwidth}{!}{\begin{tikzpicture}[
  font=\sffamily\small,
  box/.style={draw=charcoal,rounded corners=2pt,very thick,align=center,
    inner sep=7pt,minimum height=23mm},
  coord/.style={box,fill=formalblue!9,draw=formalblue},
  map/.style={box,fill=warningamber!10,draw=warningamber},
  obl/.style={box,fill=passgreen!9,draw=passgreen},
  invariant/.style={box,fill=softgray,draw=charcoal},
  open/.style={box,fill=stopred!7,draw=stopred,dashed},
  arr/.style={-{Latex[length=1.6mm,width=1.2mm]},very thick,draw=charcoal},
  darr/.style={-{Latex[length=1.6mm,width=1.2mm]},very thick,dashed,draw=stopred}
]
\node[coord,minimum width=42mm] (coordinate) at (0,0) {
  coordinate table\\
  $\epscoord_n\epscoord_m=\delta_{nm}\epscoord_n$\\
  orthogonal in source coordinates};
\node[map,minimum width=42mm] (zeta) at (5.6,0) {
  unitriangular change\\
  $\zeta\in1+J_N$\\
  incidence off-diagonal data};
\node[obl,minimum width=44mm] (oblique) at (11.5,0) {
  compiled table\\
  $q_n=\zeta\epscoord_n\zeta^{-1}$\\
  oblique, primitive, complete};

\node[invariant,minimum width=105mm] (cyclic) at (5.75,-2.75) {
  ordinary cyclic invariant is unchanged:\quad
  $\displaystyle\Tr Q(b)^r=\sum_n b_n^r$,\qquad
  $\displaystyle\det(I-zQ(b))=\prod_n(1-zb_n)$};
\node[open,minimum width=73mm] (adjoint) at (5.75,-4.65) {
  not covered by ordinary similarity invariants:\\
  mixed adjoint Gram geometry $q_p^*q_q$};

\draw[arr] (coordinate.east) -- node[above,font=\scriptsize] {conjugate} (zeta.west);
\draw[arr] (zeta.east) -- (oblique.west);
\draw[arr] (coordinate.south) -- ++(0,-0.45) -| (cyclic.north west);
\draw[arr] (oblique.south) -- ++(0,-0.45) -| (cyclic.north east);
\draw[darr] (cyclic.south) -- (adjoint.north);
\end{tikzpicture}
}
\caption{At finite cutoff the incidence compiler is a unitriangular change of
basis from the coordinate projectors. Its off-diagonal geometry is genuine,
but every ordinary cyclic trace and determinant is unchanged. Adjoint Gram
products are deliberately placed beyond the dashed boundary: they are not
ordinary cyclic invariants and form the separate forward obligation.}
\label{fig:similarity-firewall}
\end{figure}

\subsection{Why scalar Möbius inversion is not the selector}

The two uses of Möbius structure must not be conflated. In
\eqref{eq:compiler}, the incidence inverse is part of a conjugation that
orthogonalizes a complete coordinate family. The scalar arithmetic function
\(\mob(n)\) is not itself an atom indicator:
\[
\mob(2)=-1,\qquad \mob(6)=1,\qquad \mob(4)=0.
\]
It fails idempotence at \(2\), accepts the squarefree composite \(6\), and
vanishes for a reason unrelated to primality at \(4\).

The von Mangoldt function has a different limitation. Its prime-power support
appears after scalar endpoint convolution and does not define a primitive
word action. Likewise, prime-zeta inversion acts after commutative
aggregation. Neither distinguishes the source letter \(p^r\) from the
temporal word \(p,p,\ldots,p\) at the operator level.

The unfiltered compiler gives the complementary ablation. If one defines
\(A_n=q_n\) for every \(n\), every composite coordinate survives as an
idempotent. Möbius conjugation makes coordinates orthogonal; it does not say
which coordinates are atoms. That ownership belongs entirely to the cover
predicate \eqref{eq:atom-predicate}.

\begin{table}[t]
\centering
\small
\caption{Ablations separating atom selection from incidence compilation.}
\label{tab:algebraic-ablations}
\begin{tabularx}{\textwidth}{L{37mm} L{43mm} Y}
\toprule
ablation & exact outcome & ownership lesson \\
\midrule
scalar \(\mob(n)\)
& \(\mob(2)=-1,\ \mob(6)=1\)
& scalar inversion is neither an atom predicate nor an idempotent weight \\
\(\zeta\epscoord_n\) without \(\mu\)
& distinct columns need not annihilate
& the inverse is required to compile a complete orthogonal family \\
all \(q_n\), no cover filter
& every composite coordinate survives
& inversion compiles coordinates; covers select atoms \\
mutated relation with \(6\) covering \(1\)
& the compiler selects \(6\)
& selection is equivariant with the supplied grammar and proves too much \\
coordinate relabeling
& output is conjugately relabeled
& numeric prime names are not hidden in the formula \\
\bottomrule
\end{tabularx}
\end{table}

\subsection{The exact negative conclusion}

The no-go does not say that every \(q_n\) is literally diagonal in the
original Hilbert metric. It says that their ordinary cyclic character is
coordinate-like and, at finite cutoff, that the entire family is explicitly
similar to the coordinate table. Radical extensions and obliqueness cannot
alter \eqref{eq:finite-power-trace} or \eqref{eq:finite-determinant}.

This conclusion is stronger than determinant equivalence discovered only
after computation. It is a classification statement with an explicit
conjugating unit. It also points to the only datum not tested here:
\(q_p^*q_q\) depends on the chosen Hilbert geometry and need not vanish even
when \(q_pq_q=0\). That adjoint question is not used to repair the present
route.
